\cotefile{34}
\scanpage{100}
\mspage{15}
% archive sheet number: 99
% source lines 1--4
sans supposer $U$ quasi-projectif sur $K$, mais seulement de t.f. et séparé. D'autre part, les réductions \deleted{faites à la fin de} \inserted{dans} I nous permettent de conclure :

% source lines 5--9
\uline{Corollaire} 3.7 Sous les conditions (i) (ii) de 3.5, pour tt schéma \inserted{$U$} de type \inserted{fini} sur $K$, $H^i(U_{\bar K},\unsure{\mathbb Z_\ell|Z_\ell})$ est de t.f. sur $\mathbb Z_\ell$ et l'action du groupe d'inertie $\bar I$ sur \deleted{\unsure{H}} lui est \unsure{qsi|quasi} unipotente. Même \unsure{résultats|résultat} pour $H^i_!(U_{\bar K},\mathbb Z_\ell)$, quand $U$ est \deleted{\unsure{de t.f.}} \inserted{séparé} sur $K$.

% source lines 10--23, including the interlinear revisions
\uline{Remarques} 3.8. \deleted{\unsure{La démonstration que nous donnons}} \inserted{\unsure{a)|a°)} Pour l'instant, 3.5 \deleted{\unsure{n'est applicable}} \inserted{et 3.7,}} \unsure{sont applicables} tels quels que si $S$ est de caractéristique nulle, \deleted{de 3.2 montre que} en utilisant les résultats de Hironaka [\ ] et d'Artin SGA 4 $\overline{\underline{\mathrm{XIX}}}$ \ldots\ . D'autre part, \deleted{\unsure{pour cela,}} lorsqu'on \unsure{disposera|dispose} de la résolution des singularités \deleted{\unsure{sous la forme où}} pour les schémas de type fini sur un corps $k_0$ \inserted{(en supposant \unsure{s. dém. algébrique})}, alors, utilisant le th. de pureté relatif SGA 4 $\overline{\underline{\mathrm{XV}}}$ \ldots\ , on voit \deleted{que le \unsure{même rais.},} \inserted{aisément (en procédant à peu de choses près comme pour 3.5), que le} th. de monodromie est valable \deleted{pour les schémas,} \deleted{\unsure{sous}} sous les formes 3.5 et 3.7, lorsque $S$ est \deleted{local} hensélisé \inserted{\unsure{--|ou} hensélisé strict} d'un schéma de type fini sur \deleted{un corps} $k_0$.

% source lines 24--26
\deleted{\unsure{Par a.}} b) Lorsque $n=1$, \unsure{alors|on a} grâce à Abhyankar [\ ] la condition (ii) de 3.5. est tjs satisfaite. On ne sait apparemment s'il en est de $=$ de la
\scanpage{101}
\mspage{16}
% archive sheet number: 100
% source lines 1--9
condition (i). Cependant, la démonstration que nous donnerons \inserted{au § 5} de (3.2) montre que \uline{si le th. de pureté est satisfait pour les} $(X,Y^{(i)})$ \uline{\unsure{qui ont} dimension} $\leq i$, \uline{alors la conclusion sur} $H^i(U_{\bar K},\unsure{\mathbb Z_\ell|Z_\ell})^P$ \uline{reste valable}. Or \deleted{\unsure{la}} cette condition est toujours satisfaite pour $i=1$, \unsure{où|car} elle résulte du th. de pureté \unsure{en dimension} $2$ de \inserted{Nagata--}Zariski (SGA 2 \ldots), \deleted{Nous} cf plus bas \ldots\ . Le raisonnement de 3.5 a donne donc \deleted{\unsure{également}} ceci :

% source lines 10--13
\deleted{\uline{\unsure{Corollaire}} \unsure{3.9}.} Soit $X_K$ une \deleted{\unsure{schéma lisse de}} \inserted{courbe algébrique} \deleted{lisse de dim. $\leq 1$ sur $K$.} lisse sur $K$, alors l'action du groupe d'inertie $\bar I$ sur $H^1(X_{\bar K},\mathbb Z_\ell)$ est \unsure{qsi.|qu.} unipotente \uline{d'échelon 2}.

% source line 14
Grâce à la réduction II 5.4, on en conclut

% source lines 15--19
\uline{Corollaire} 3.9. Soit $X_K$ un schéma \deleted{normal} de type fini sur $K$. Alors $E=H^1(X_{\bar K},\mathbb Z_\ell)$ est de t.f., \inserted{et} l'action du groupe d'inertie sur $E$ est \unsure{qsi|quasi} unipotente d'échelon 3, \deleted{\unsure{et}} et $=$ qsi unipotente d'échelon 2 si $X$ est normal.
\scanpage{102}
\mspage{17}
% archive sheet number: 101
\margin{$\overline{\underline{\mathrm{III}}}$}
% source heading
4. \uline{Calculs \inserted{cohomologiques} locaux liés au théorème de pureté}.

% margin note slanted
\margin{devrait \unsure{remonter|commencer} au n° 3}
% source lines 1--7
\uline{Définition} 4.1. a) Soit $X$ un schéma régulier. On dit que $X$ \uline{satisfait au théorème de pureté cohomologique} si, pour tout \deleted{\unsure{morphisme}} $X'$ étale sur $X$, et tout \deleted{de} sous-schéma régulier $Y$ de $X'$, \deleted{purement de} tout sous-schéma régulier $Z$ de $Y$ purement de codimension 1, \deleted{\unsure{on a}} \deleted{$\underline H^i_Z($} et tout entier $n$ premier aux \deleted{\unsure{la}} caractéristiques résiduelles de $Z$, on a (si $\Lambda=\unsure{\mathbb Z/n\mathbb Z|Z/nZ}$)

\[
\underline H^i_Z(\Lambda_Y)=0\qquad\text{pour }i\ne 2
\tag*{(4.1.1)}
\]

% source lines 9--10
ou ce qui revient au même, si $f:U=Y-Z\unsure{{-}|\longrightarrow}Y$ est l'inclusion canonique

\[
R^i f_*(\Lambda_U)=0\qquad\text{pour }i\ne 0,1.
\tag*{(4.1.2)}
\]

% source lines 12--14
\deleted{\unsure{R}} \inserted{b) Si $r$ est un entier,} on dit que $X$ satisfait au théorème de pureté cohomologique en dim $\leq r$ si \deleted{\unsure{4.1.2}} (4.1.2) est vrai pour $i\leq r$.

% source lines 15--20
\uline{Remarques} 4.2. a) La condition énoncée \deleted{dans 4.1 a} \deleted{\unsure{pour}} est \deleted{\unsure{j'expr.}} évidemment locale pour la top. (ét). Elle peut s'énoncer \inserted{aussi} en disant que pour tt \uline{localisé strict} $X'$ de $X$, tt sous-schéma régulier $Y$ de $X'$, et tt diviseur régulier $Z$ de $Y$, \inserted{et tt $n$ premier à la car. rés. de $X'$,} on a (4.1.1) ou, ce qui revient au $\unsure{\overline{m}|\overline{\sim}}$

\[
H^i(Z-Y,\mathbb Z/n\mathbb Z)=0\qquad\text{si }i\ne 0,1.
\]

% source lines 22--24
(resp. et $i\leq r$). \deleted{\illegible{3}}

b) Il est évident \inserted{(simple définition)} que \deleted{\unsure{si l'on}} $X$ satisfait au th. de pureté \inserted{coh.} en dimensions $\leq 1$.
\scanpage{103}
\mspage{18}
% archive sheet number: 102
\margin{\deleted{\unsure{SGA}}}
% source lines 1--4
c) Rappelons qu'on conjecture que tt schéma régulier satisfait au th. de pureté cohomologique (SGA 4 $\overline{\underline{\mathrm{XVI}}}$ 3.10 b). Cette propriété est démontrée pour l'instant dans les cas suivants :

% source lines 5--7
1°) $X$ est lisse sur un corps parfait (plus gén., les relations (4.1.1) ou (4.1.2) sont vraies si $Y,Z$ sont lisses sur un corps quelconque), cf. (SGA 4 $\overline{\underline{\mathrm{XVI}}}$ 3.7, 3.9).

% source lines 8--9
2°) $X$ est excellent de caractéristique nulle (SGA 4 XIX 3.2). \deleted{Plus généralement}

% source lines 10--16, interlinear revision of pour aucun
\deleted{Mais pour} Plus généralement, \inserted{on prouve dans loc. cit. que} si $X$ est excellent de car. donnée, \deleted{\unsure{et}} et si on dispose de la résolution des singularités des schémas excellents en cette caractéristique, alors $X$ satisfait au th. de pureté cohomologique. Par contre, \deleted{\unsure{dans ce cas}} pour \deleted{un} \inserted{aucun} schéma régulier d'inégales caractéristiques et de dimension $\geq 2$, le théorème de pureté n'est prouvé à l'heure actuelle.

% interlinear addition below the preceding paragraph
\inserted{d) \unsure{Les|La} condition de 4.1. est stable par localisation \uuline{étale} et par passage à un sous-schéma régulier.}

% source lines 17--20
\uline{Proposition} 4.2. Supposons que le th. de pureté \deleted{soit} cohomologique soit vrai pour $X$, et soit $Y$ un sous-schéma fermé régulier de $X$, purement de codimension $d$. Alors on a, si $\Lambda=\unsure{\mathbb Z/n\mathbb Z|Z/nZ}$, ($n$ premier aux car. rés de $X$)

\[
\underline H^i_Y(\Lambda_X)=
\deleted{\left\{\begin{aligned}\unsure{0}\quad\unsure{i\ne 2d}\\ \mu_n^{\otimes(-d)}\quad\text{si}\end{aligned}\right.}
\unsure{0}\qquad\text{si }i\ne 2d,
\tag*{(4.2.1)}
\]

% source line 22
et on a un isom. canonique

\[
\underline H^{2d}_Y(\Lambda_X)\simeq(\mu_n^{\otimes(-d)})_Y\quad\deleted{\unsure{\text{si}}}
\tag*{(4.2.2)}
\]

(ou encore

\[
\underline H^{2d}_Y(\mu_n^{\otimes d})\simeq\deleted{\unsure{\mathbb Z/n\mathbb Z}}\Lambda_Y\qquad).
\tag*{(4.2.2 bis)}
\]
\scanpage{104}
\mspage{19}
% archive sheet number: 103
% source lines 1--6, top additions and cancellations
\uuline{Démonstration} \deleted{Pour (4.2.1), la question étant locale,} \deleted{la \unsure{pureté} \unsure{par|pour}} \inserted{Prouvons d'abord (4.2.1) et le fait que} \deleted{\unsure{le}} $\underline H^{2d}_Y(\Lambda_X)$ est loc. isom. à $\Lambda_Y$. \deleted{\unsure{raisonnons sur $d$,}} Le cas $d=0$ est trivial, le cas $d=1$ résulte de la théorie de Kummer et du \unsure{Lemme|lemme} de Abhyankar absolu, \deleted{\unsure{déduisant}} \inserted{qui fournit $\underline H^2_Y(\mu_{nX})\simeq\Lambda_Y$} (\deleted{SGA 4} cf SGA 4 $\overline{\underline{\mathrm{XIX}}}$ 1.3 ou SGA 1 $\overline{\underline{\mathrm{XII}}}$ \ldots).

% source lines 7--11
Pour $d>1$, on procède par récurrence sur $d$. L'assertion étant locale, on \deleted{\unsure{se ramène}} est ramené au cas où il existe un \uline{diviseur régulier} $Z$ dans $X$ tel que \deleted{\unsure{$Z$}} $Y\subset Z$. \deleted{On utilise} \inserted{Considérons} La suite spectrale de transitivité

\[
\underline H^*_Y(\Lambda_X)\Longleftarrow E^{pq}_2=\underline H^p_Y(\underline H^q_Z(\Lambda_X)),
\]

% source lines 13--17
\deleted{et l'hyp. de réc.} alors \deleted{\unsure{$E^p$}} $\underline H^q_Z(\Lambda_X)$ est nul si $q\ne 2$, et loc. isom. à $\Lambda_Y$ si $\unsure{q|p}=1$, \unsure{donc|dans} l'hyp. de réc. nous prouve que $E^{pq}_2=0$ sauf \deleted{\unsure{p.}} si \deleted{$p=d-1$,} $p=2(d-1)$, $q=2$, et que pour ces valeurs de $p,q$, $E^{pq}_2$ est loc. isom. à $\Lambda_Z$, d'où la conclusion.

% source lines 18--22
Pour la définition d'un isom. canonique (4.2.1) ou (4.2.1 bis), le cas $d=1$ est traité dans (SGA 4 XIX 1.3), \deleted{le cas $d$} Pour le cas $d$ quelconque, il suffit d'utiliser la construction de (SGA 4 $\overline{\underline{\mathrm{XIX}}}$ 3.4) ou de (SGA 4 $\overline{\underline{\mathrm{XVIII}}}$\quad).
\scanpage{105}
\mspage{20}
% archive sheet number: 104
% source lines 1--6
\uline{Corollaire} 4.3 Soit $X$ un \deleted{\unsure{sous-}}schéma \deleted{fermé} régulier \deleted{\unsure{au th. de pureté}} \inserted{satisfaisant à la} pureté coh. (4.1), $Y$ un sous-schéma fermé régulier purement de codim. $d$, $n$ un entier premier aux caractéristiques résiduelles de $X$, $\Lambda=\unsure{\mathbb Z/n\mathbb Z|Z/nZ}$, $F$ un faisceau loc. constant sur $X$ annulé par $n$.
\inserted{Pour tout faisceau de $\Lambda$-modules $G$ sur $X$ et tt $d\in\mathbb Z$, on pose $G(d)=G\otimes\mu_n^{\otimes(d)}$.}
\deleted{Alors} on a une suite exacte canonique

\[
\cdots\unsure{{-}|\longrightarrow}H^{i-2d}(Y,F\deleted{\unsure{\otimes\mu_n^{\otimes(-d)}}}(-d))
\xrightarrow{\alpha}H^i(X,F)\xrightarrow{\beta}H^i(X-Y,F)
\xrightarrow{\gamma}H^{i-2d+1}(Y,F\deleted{\unsure{\otimes\mu_n^{\otimes(-d)}}}(-d))\longrightarrow\cdots
\tag*{(4.3.1)}
\]

% source lines 8--9
où $\beta$ est l'homomorphisme de restriction, et où l'homomorphisme composé

\[
H^{i-2}(X,F\deleted{\unsure{\otimes\mu_n^{\otimes d}}}(d))
\xrightarrow{\text{restr.}}H^i(Y,F\deleted{\unsure{\otimes\mu_n^{\otimes(-d)}}}(\unsure{-d|d}))
\xrightarrow{\alpha}H^i(X,F)
\]

% source lines 11--12
est le cup-produit avec une classe de cohomologie \deleted{bien déterminée} \inserted{canonique}

\[
\gamma(Y)\in\deleted{\unsure{\text{de}}}H^{2d}(X,\deleted{\unsure{\mu_X}}\Lambda_X(d)).
\tag*{(4.3.2)}
\]

% source lines 14--17 and additions
\deleted{Quand $Y$ est \unsure{de}} Lorsque $d=1$, \inserted{\unsure{de sorte que} $Y$ est un \uline{diviseur} sur $X$,} cette classe est aussi la \inserted{première} classe de Chern de $\underline L=\unsure{\mathscr O_X|\mathcal O_X}(Y)$, i.e. l'image \inserted{de la classe de $\underline L$ dans $\operatorname{Pic}(X)\simeq H^1(X,\unsure{\mathbb G_{mX}|G_{mX}})$} par l'homomorphisme \deleted{\unsure{de}} de cobord associé à la suite exacte de Kummer

\[
0\longrightarrow\mu_{nX}\longrightarrow\mathbb G_{mX}\overset{\lambda\mapsto\lambda^n}{\unsure{{-}|\longrightarrow}}\mathbb G_{mX}\longrightarrow 0.
\]

% source lines 19--20
\uline{Dém.} La suite exacte (4.3.1) n'est autre essentiellement que la suite exacte de cohomologie relative

\[
\cdots\longrightarrow H^i_Y(X,F)\unsure{{-}|\longrightarrow}H^i(X,F)\unsure{{-}|\longrightarrow}H^i(X-Y,F)\longrightarrow\cdots,
\]

% source line 22
compte tenu des isomorphismes

\[
H^i_Y(X,F)\simeq H^{i-2d}(Y,F\deleted{\unsure{\otimes\mu_n^{\otimes(-d)}}}(-d))
\]

% source line 24
déduits de la suite spectrale de localisation

\[
H^{\unsure{*|n}}_Y(X,F)\Longleftarrow\unsure{E|F}^{pq}_2=H^p(Y,\underline H^q_Y(F)),
\]
\scanpage{106}
\mspage{21}
% archive sheet number: 105
% source lines 1--2
et du calcul 4.2 des $\underline H^q_Y(F)$, impliquant que $E^{pq}_2=0$ si $q\ne 2d$, d'où $H^i_Y(X,F)\simeq E^{i-2d,2d}_2$.

% source lines 3--4
Faisant \deleted{\unsure{$F$}} $F=\mu_n^{\otimes d}$ et $i=2d$ dans (4.3.1), on trouve un homomorphisme

\[
H^{2d}_Y(X,\deleted{\unsure{\mu_n^{\otimes d}}}\Lambda(d))\simeq H^0(Y,\Lambda)\longrightarrow H^{2d}(X,\deleted{\unsure{\mu_n^{\otimes d}}}\Lambda(d))
\]

% source lines 6--11
et l'image de la section \deleted{1 \unsure{dans le premier}} \inserted{\unsure{constante} 1 de $\Lambda_Y$ sur $Y$} nous donne un élément canonique (4.3.2). Les assertions \deleted{\unsure{les}} de compatibilité faites \inserted{dans 4.3} sur cet élément \deleted{dans} se prouvent par des arguments standard, pour lesquels nous renvoyons à SGA 5 $\overline{\underline{\mathrm{IV}}}$.

% source lines 12--14
\uline{Remarque} 4.3.3. On a \inserted{\unsure{seulement}} utilisé l'hyp. de pureté faite pour \deleted{conclure} avoir \deleted{\unsure{les}} les \unsure{valeurs} de $\underline H^i_Y(\Lambda_X)$ \ldots\ .
\scanpage{107}
\mspage{22}
% archive sheet number: 106
% cancelled opening, source lines 1--3
\deleted{4.3. \unsure{Soit} Soit $X$ un schéma régulier, et soit}

\deleted{Rappelons \unsure{pour mémoire} une forme du th. de Abhyankar}

% source lines 4--8
\unsure{4.4.|4.3.} Soit $X$ un schéma régulier, $(\unsure{D_i|Z_i})_{1\leq i\leq r}$ une famille de \deleted{\unsure{sché.} diviseurs premiers rég.} sous-schémas réguliers irréductibles dans $X$, tels que le diviseur $D=\sum D_i$ soit un diviseur à croisements normaux (SGA 5 I \ldots).

% source lines 9--13
Soit $U=X-\operatorname{supp}D=X-\bigcup D_i$, \inserted{$f:U\longrightarrow X$ l'inclusion,} $n$ un entier premier aux caractéristiques résiduelles de $X$ en les points de \deleted{\unsure{$Z$}} $\operatorname{supp}D$, $\Lambda=\unsure{\mathbb Z/n\mathbb Z|Z/nZ}$. On se propose de calculer les faisceaux $R^i f_*(\Lambda_U)$ sur $X$. \deleted{$\unsure{R^i}f_*(\Lambda_U)$}

% source lines 14--15
\uline{Proposition} 4.5. \deleted{Avec} Les notations \inserted{étant celles} de 4.4, \deleted{et \unsure{supposons} que $X$ satisfait au th. de pureté}

(i) \deleted{$R^1 f_*(\Lambda_U)$ est canoniquement isom.} \inserted{On a un isom. canonique}

\[
R^1f_*(\Lambda_U)\simeq\deleted{\unsure{\text{le}}}\coprod_{1\leq i\leq r}(\mu_n^{\unsure{\otimes(-1|\otimes(-1)}})_{D_i}\simeq\coprod_{1\leq i\leq r}\Lambda_{D_i}(-1)
\tag*{(4.\deleted{4}5.1)}
\]

% boxed left-margin insertion keyed to (ii)
(ii) \inserted{Supposons que $X$ satisfait au th. de pureté cohomologique (4.1)} \inserted{Pour tout entier $p$,} L'homomorphisme \inserted{suivant,} déduit des cup-produits est un \deleted{\unsure{produit}} isomorphisme

\[
\bigwedge^p R^1 f_*(\Lambda_U)\xrightarrow{\sim}R^p f_*(\Lambda_U).
\tag*{(4.\deleted{4}5.2)}
\]

(iii). Par suite, \inserted{sous l'hyp. de (ii),} on a un isomorphisme canonique

\[
\begin{aligned}
R^p f_*(\Lambda_U)&\simeq\coprod_{1\leq i_1<i_2<\cdots<i_p\leq r}(\deleted{\unsure{\mu_n^{\otimes(-p)}}}\Lambda_X(-p))_{D_{i_1\cdots i_p}},\\
&\simeq\{\deleted{\unsure{\mu_n^{\otimes(-p)}}}\Lambda_X(-p)\}\otimes\coprod_{1\leq i_1<\cdots<i_p\leq r}\Lambda_{D_{i_1\cdots i_p}},
\end{aligned}
\tag*{(4.5.3)}
\]

% cancelled trial coefficient above the right part of (4.5.3)
\margin{\deleted{$\unsure{(\mu_n^{\otimes p})_X}$}}
% lower-left margin note
\margin{écrire $\Lambda_X(-p$}
% source last line, after a cancelled beginning
\deleted{\unsure{si} \illegible{3} \unsure{posé}} où on a posé $D_{i_1\cdots i_p}=D_{i_1}\unsure{\cap|\wedge}D_{i_2}\cap\cdots\cap D_{i_p}$.
\scanpage{108}
\mspage{23}
% archive sheet number: 107
% source lines 1--3
\uline{Démonstration}. Soit $U_i=X-D_i$, $f_i:U_i\unsure{{-}|\longrightarrow}X$ l'inclusion, \deleted{de sorte qu'il \unsure{donne un}} \inserted{de sorte que} $f:U\unsure{{-}|\longrightarrow}X$ se factorise en

\diagram{\begin{tikzcd}[ampersand replacement=\&]
U \arrow[r,"g_i"] \arrow[rr,bend left=35,"f"] \& U_i \arrow[r,"f_i"] \& X.
\end{tikzcd}}

% source line 5
On en conclut un hom canonique

\[
\deleted{\unsure{(a)}}\qquad R^1 f_{i*}(\Lambda_{U_i})\unsure{{-}|\longrightarrow}R^1 f_*(\Lambda_U)
\]

% source line 7 and its correction
D'autre part, \deleted{(4.2.2) \unsure{nous dit et th. canonique}} \inserted{on a un isom (4.2.2)}

\[
R^1 f_{i*}(\Lambda_{U_i})\simeq\underline H^2_{D_i}(\Lambda_X)\simeq\deleted{\unsure{(\mu_n^{\otimes-1})_{D_i}}}\Lambda_{D_i}(-1),
\]

d'où un homomorphisme

\[
\unsure{\varphi_i|\varphi}:\deleted{\unsure{(\mu_n^{\otimes-1})_{D_i}}}\Lambda_{D_i}(-1)\longrightarrow R^1 f_*(\Lambda_U),
\]

et par suite un hom

\[
\deleted{\unsure{\text{avec}}}\qquad\varphi:\coprod_{1\leq i\leq r}\deleted{\unsure{(\mu_n^{\otimes-1})_{D_i}}}\Lambda_{D_i}(-1)\longrightarrow R^1 f_*(\Lambda_U).
\]

% source lines 14--15
Je dis que cet homomorphisme est un isomorphisme,

% The following block through the final words is crossed out by long diagonal strokes.
\deleted{Pour ceci, notons que $U=\bigcup_{1\leq i\leq r}U_i$, alors on a la suite spectrale de Leray}

\[
\deleted{R^{\unsure{*|p}}f_*(\Lambda_U)\Longleftarrow E^{pq}_2=H^p((i_0\cdots i_p)\longrightarrow R^q f_{i_0\cdots i_p*}(\Lambda_{U_{i_0\cdots i_p}})\unsure{)})}
\]

\deleted{où $U_{i_0\cdots i_p}=U_{i_0}\cap\cdots\cap U_{i_p}$, et où \deleted{\unsure{$f$.}} $f_{i_0\cdots i_p}:U_{i_0\cdots i_p}\longrightarrow X$ est l'inclusion canonique.}

\deleted{On a évidemment $R^0 f_{i_0\cdots i_p*}(\Lambda_{U_{i_0\cdots i_p}})\simeq\Lambda_X$, donc $E^{p0}_2=0$ pour tout $p\ne 0$, d'où résulte que la suite spectrale donne un isomorphisme}

\[
\deleted{\deleted{\unsure{E_2}}\quad R^1 f_*(\Lambda_U)\xrightarrow{\sim}E^{01}_2\subset\coprod_{1\leq i\leq r}R^1 f_{i*}(\Lambda_{U_i})}
\tag*{\deleted{(***)}}
\]

\deleted{Or on constate tout de suite que le}

% small drawing in the lower-left margin
\figure{\begin{tikzpicture}[x=0.6cm,y=0.6cm]
\draw (-0.65,0)--(3.7,0);
\draw (0,-0.6)--(0,3.4);
\node at (0,0) {$\times$};
\node at (0,0.48) {$\times$};
\node at (0.9,0.52) {$\times$};
\node at (1.7,0.57) {$\times$};
\node at (2.55,0.68) {$\times$};
\end{tikzpicture}}
\scanpage{109}
\mspage{24}
% archive sheet number: 108
% Lines 1--12: continuation of the large diagonally cancelled block on scan 108.
\deleted{composé de $\deleted{\unsure{\varphi|\varphi_i}}\inserted{\varphi_i}$ avec l'isom précédent est l'inclusion canonique de $R^1f_{i*}(\Lambda_{U_i})$ dans la somme au dernier \unsure{mb|membre} de (***). Donc on en conclut \deleted{\unsure{que}} \deleted{\unsure{\(\varphi\)}} aussitôt que \deleted{\illegible{1}} $E_2^{01}$ est égal au dernier membre de (***), et que $\unsure{\varphi|\varphi_i}$ est un isomorphisme. Cela prouve (i). \deleted{Pour} Supposons maintenant que $X$ satisfait au th. de pureté, et prouvons (ii). On \deleted{\unsure{peut supposer}} \inserted{\unsure{se ramène à}} vérifier l'assertion \deleted{\unsure{fibre par fibre}} \unsure{géométrique}, donc on est ramené \inserted{au cas} $X$ strictement local, et \deleted{\unsure{se trouve à pouvoir}} prouver alors que}

\[
\deleted{\bigwedge^p H^1(U,\Lambda)\longrightarrow H^p(U,\Lambda)}
\]

\deleted{est un isomorphisme.}

% Lines 13--16: live replacement begins below the cancellation boundary.
en d'autres termes que l'hom. canonique

\[
\coprod_{1\leq i\leq r}(\Lambda)_{D_i}\longrightarrow R^1f_*(\deleted{\unsure{\Lambda_X|\Lambda_U}}\unsure{\mu_X|\mu_{nX}}^{\deleted{\unsure{-1|\otimes(-1)}}})
\tag*{(\unsure{4.5.4|4.5.1})}
\]

% Lines 17--23.
est un isomorphisme. On \deleted{\unsure{peut se}} est ramené pour ceci, \deleted{\unsure{aux propriétés}} \inserted{ainsi que pour (ii),} à vérifier l'assertion fibre par fibre, ce qui nous ramène au cas où $X$ est str. local, et à justifier alors que l'hom. \unsure{correspondant à ces}

\[
\deleted{\unsure{\coprod_{1\leq i\leq r}|\coprod_{i\leq r}}\unsure{\Lambda|\mu_n}}\Lambda^r\longrightarrow H^1(U,\mu_n\deleted{\unsure{\Lambda}})
\]

est un isom, ainsi que l'hom.

\[
\bigwedge^d H^1(U,\deleted{\unsure{\mu_n}}\Lambda)\longrightarrow H^d(U,\Lambda)\deleted{\unsure{*|\Lambda}}
\]

dans le cas où $X$ satisfait au th. de pureté cohomologique.
\scanpage{110}
\mspage{25}
% archive sheet number: 109
% Lines 1--6.
La première assertion résulte \inserted{aussitôt} de la détermination du ``groupe fondamental modéré'' de $U$, v. plus bas (\quad). Pour la deuxième, on procède par récurrence sur $r$, le cas $r=0$ étant trivial, et le cas $r=1$ \deleted{\unsure{est}} \unsure{étant} est contenu dans 4.2 (où on fait $d=1$). \deleted{\unsure{Si}} Si $r\geq2$, on \deleted{écrit} considère

\[
\begin{aligned}
U'&=X-\deleted{\unsure{\sum|D_i}}\bigcup_{1\leq i\leq r-1}D_i\\
D'_r&=\deleted{\unsure{X\cap D_r}}U'\cap D_r
\end{aligned}
\]

de sorte que

\[
\begin{aligned}
U&=U'-D'_r\ ;\ \deleted{\unsure{D'_r=D_r-\sum_{1\leq i\leq r-1}D_i\cap D_r}}\\
D'_r&=D_r-\operatorname{supp}\Delta_r,\qquad \Delta_r=\sum_{1\leq i\leq r-1}D_i\cap D_r.
\end{aligned}
\]

\margin{\deleted{\unsure{la fibrée|la fibre}}}

On a la suite \unsure{ex.act|ex.acte} de cohomologie relative (4.3.1) pour $D'_r$ dans $U'$

% From this display to the bottom, the entire block is diagonally cancelled.
\[
\deleted{
\cdots\longrightarrow H^{i-2}_{\deleted{\unsure{D'_r}}}(D'_r,\deleted{\unsure{\Lambda}}\mu_n^{\unsure{\otimes-1|\otimes1}})\xrightarrow{\alpha_i}H^i(U',\Lambda)\longrightarrow H^i(U,\Lambda)
\longrightarrow H^{i-1}_{\deleted{D'_r}}(D'_r,\deleted{\unsure{\Lambda}}\mu_n^{\unsure{\otimes-1|\otimes1}})\longrightarrow\cdots.
}
\tag*{\deleted{(4.5.5)}}
\]

\[
\deleted{\inserted{H^{i-2}(D'_r,\Lambda)\otimes\mu_n(\unsure{k})^{-1}}}
\]

\deleted{D'autre part \deleted{\unsure{on a}} \unsure{la} suite spectrale}

\[
\deleted{H^{\unsure{*|\bullet}}_{D'_r}(U,\Lambda_U)\Longleftarrow \inserted{E_2^{pq}=}H^p(\unsure{D_r|D'_r},\underline{H}^{q}_{D'_r}(\deleted{\unsure{\mu_n}}\Lambda_{U'}))}
\]

\deleted{et la détermination 4.2 des $\underline{H}^q_{D'_r}(\Lambda_{U'})$ donne}

\[
\deleted{H^i_{D'_r}(U',\Lambda_{U'})\simeq H^{i-2}(D'_r,\mu_{nD'_r}^{\otimes-1})\simeq H^{i-2}(D'_r,\Lambda)\otimes\mu_n(k)^{-1}}
\]

\deleted{où $k$ est le corps résiduel de $X$.}

% A rewritten paragraph is inserted above the spectral-sequence paragraph; its continuation runs down the left edge.
\deleted{\inserted{L'hypothèse de récurrence, appliquée à $D'_r$ et les $r-1$ div. réguliers $D'_i$ $(1\leq i\leq r-1)$ sur $D'_r$}}

\deleted{\deleted{\unsure{D'autre part}, l'hypothèse de récurrence appliquée à $D'_r$ et aux \inserted{$r-1$ diviseurs réguliers} $(D_i\cap D_r)$ $(1\leq i\leq r-1)$ sur $D'_r$} implique:}

\[
\deleted{H^{\deleted{\unsure{*}}i}(D'_r,\Lambda)\simeq\bigwedge^{\deleted{\unsure{*}}i}H^1(D'_r,\Lambda)\simeq\left(\bigwedge^{\deleted{\unsure{*}}i}(\Lambda^{r-1})\right)\otimes\mu_n^{\otimes-i}(k)}
\]

\deleted{et de même l'hypothèse de récurrence appliquée à $X$ et aux $D_i$, $(1\leq i\leq r-1)$ donne}

\[
\deleted{H^i(U',\Lambda)\simeq\bigwedge^iH^1(U',\Lambda)\simeq\left(\bigwedge^i(\Lambda^{r-1})\right)\otimes\mu_n^{\otimes-i}(k).}
\]

\deleted{\deleted{Il faut exprimer \unsure{maintenant} l'hom de \unsure{(***)}}}

\[
\deleted{\deleted{H^i_{D'_r}(U,\Lambda)\simeq H^{i-2}(D'_r,\Lambda)\otimes\mu_n^{\unsure{\otimes1|\otimes-1}}(k)\simeq\bigwedge^{i-2}(\Lambda^{r-1})\otimes\mu_n^{\deleted{\unsure{\otimes(-i+1)|\otimes(-i)}}\inserted{\unsure{-(i-2)-1}}}(k)}}
\]

% A bent annotation arrow leads from the preceding first term to this added formula.
\[
\deleted{\inserted{H^i(U',\Lambda)\simeq\bigwedge^i(\Lambda^{r-1})\otimes\mu_n^{\otimes-i}(k)}}
\]
\scanpage{111}
\mspage{26}
% archive sheet number: 110
% Line 1; the last arrow reaches the right edge of the sheet.
\[
\cdots\longrightarrow H^{i-2}(D'_r,\Lambda(-1))\xrightarrow{\alpha_i}H^i(U',\Lambda)\longrightarrow H^i(U,\Lambda)\longrightarrow H^{i-1}(D'_r,\Lambda(-1))\longrightarrow
\tag*{(4.5.5)}
\]

% Lines 2--5.
De plus, \deleted{$\unsure{H^*(U)|H^*\Lambda}$}, les termes $H^{\unsure{*|\alpha}}(U',\Lambda)$ et $H^*(D'_r,\Lambda(-1))$ de cette suite exacte se calculent grâce à l'hypothèse de récurrence, et ils sont resp. isom. à $\bigwedge^*H^1(U',\Lambda)$ et $(\bigwedge^*H^1(D'_r,\Lambda))(-1)$. Je dis que
% Remaining area has no further recto text; faint reversed marks show through the sheet.
\scanpage{112}
\mspage{27}
% archive sheet number: 111
\deleted{Notons que}

\[
\deleted{\inserted{H^i}\quad H^{i-2}(D'_r,\Lambda)\otimes\mu_n(k)^{\otimes-1}\longrightarrow H^i}
\]

\deleted{$\unsure{H^i|\Lambda}$} Je dis que l'hom. $\alpha_i$ de (4.5.5) est \uline{nul}. En effet, les calculs explicites précédents (liés à l'hyp. de récurrence) prouvent que

\[
H^{\unsure{*|\alpha}}(U',\deleted{\unsure{E}}F)\unsure{{-}|\longrightarrow}H^*(D'_r,F)
\]

est \deleted{\unsure{surj.}} \inserted{un isom.} si $F$ est loc. constant et annulé par $n$, p.ex. si $F=\mu_n^{\otimes-1}$. Il suffit donc de prouver que le composé

\[
H^{i-2}(U',\deleted{\unsure{\mu_n^{\otimes-1}}}\Lambda(-1))\longrightarrow H^{i-2}(D'_r,\deleted{\unsure{\mu_n}}\Lambda(-1))\xrightarrow{\alpha_i}H^i(U',\Lambda)
\]

est nul. Or en vertu de 4.3, \deleted{\unsure{c'est le}} composé est le cup-produit par la classe \deleted{de} $\gamma_{U'}(D'_r)\inserted{=c^1(\underline{\unsure{\mathcal O|O}}_{U'}(D'_r))}$, \deleted{\unsure{or|si} $D'_r$ est un diviseur \unsure{p}} cette classe est induite par la classe \deleted{\unsure{de} $\unsure{D_r}$} $\gamma_X(D_r)$, qui est nulle puisque $X$ est str. local (ou \unsure{parce} que $D_r$ est principal). Donc la suite exacte (4.5.5) se décompose en suites exactes courtes

\[
0\longrightarrow H^i(U',\Lambda)\longrightarrow H^i(U,\Lambda)\longrightarrow H^{i-1}(D'_r,\Lambda)\deleted{\unsure{\otimes\mu_n(k)^{-1}}}(-1)\longrightarrow0.
\tag*{(4.5.6)}
\]

% Cancelled annotations below the first and third cohomology terms of (4.5.6).
\[
\deleted{\underset{\left(\bigwedge^i(\Lambda^{r-1})\right)\otimes\mu_n(k)^{\otimes-i}}{\wr}\qquad\underset{\bigwedge^{i-1}}{\wr}}
\]

% The final paragraph is struck through and enclosed by cancellation lines.
\deleted{D'autre part, la décomposition de \deleted{\unsure{\(\Lambda\) de la fibre}} $\Lambda^r=\coprod_{1\leq i\leq\unsure{r|r-1}}\Lambda_{D_i}$ suivant \deleted{$\unsure{\Lambda}$} $\Lambda^{r-1}\times\Lambda$, où le dernier facteur $\Lambda$ est $\Lambda_{D_r}$, donne}
\scanpage{113}
\mspage{28}
% archive sheet number: 112
\deleted{une suite exacte canonique}

\[
\deleted{0\longrightarrow\bigwedge^i\Lambda^{r-1}\otimes}
\]

\[
\begin{aligned}
H^1(U,\Lambda)&\simeq\coprod_{1\leq i\leq r}H^1(U_i,\Lambda)\simeq\coprod_{1\leq i\leq r-1}H^1(U_i,\Lambda)\times H^1(U_r,\Lambda)\\
&\simeq H^1(U',\Lambda)\times\mu_n(X)^{-1}
\end{aligned}
\]

donc une \deleted{suite exacte} \inserted{décomposition}

\[
\bigwedge^iH^1(U,\Lambda)\simeq\bigwedge^iH^1(U',\Lambda)\times\left(\bigwedge^{i-1}H^1(U',\Lambda)\otimes\mu_n(X)^{-1}\right)
\]

\[
\deleted{0\longrightarrow\bigwedge^iH^1(U',\Lambda)}
\tag*{\deleted{(4.5.7)}}
\]

\deleted{\unsure{qui fournit}} qui définit une suite exacte

\[
0\longrightarrow\bigwedge^iH^1(U',\Lambda)\longrightarrow\bigwedge^iH^1(U,\Lambda)\longrightarrow\left(\bigwedge^{i-1}H^1(U',\Lambda)\right)(-1)\longrightarrow0
\tag*{(4.5.7)}
\]

On \deleted{\unsure{obtient}} \unsure{construit} \inserted{par cup-produit} des flèches \deleted{\unsure{dans|can.}} des termes de (4.5.7) dans les termes correspondants de (4.5.6). Nous admettrons qu'on obtient ainsi un hom. de suites exactes (c'est trivial pour le premier carré, et semble demander une vérification pour le deuxième). L'\deleted{\unsure{autre|actuelle}} hypothèse de récurrence (appliquée à \deleted{$\unsure{U'=X-\bigcup D_i}$ \unsure{resp. à} $D'_r$} $U'$ resp. à $D'_r$) implique que les flèches verticales extrêmes sont \deleted{\unsure{des isom}} \inserted{bijectives}, donc aussi la flèche verticale médiane, ce qui achève la démonstration.
\scanpage{114}
\mspage{29}
% archive sheet number: 113
4.6. Supposons, \deleted{\unsure{que dans}} dans les conditions de 4.4, que l'on ait pour tout $1\leq i\leq r$

\[
\unsure{D_i|Z_i}=V(a_i),\qquad\text{où }a_i\in\Gamma(\underline{\unsure{O_X|\mathcal O_X}}).
\]

Soit $m$ un entier premier aux car. résiduelles de $X$ \deleted{\unsure{\(\sum D_i\)}} et considérons le schéma

\[
X'=\operatorname{Spec}\unsure{\mathcal O_X|\mathscr O_X}[T_1,\ldots,T_r]/(T_1^m-a_1,\ldots,T_r^m-a_r).
\]

Il est \unsure{alors|donc} bien connu que $X'$ est un schéma régulier, et que \deleted{\unsure{les}} si pour tout $1\leq i\leq r$, \deleted{\unsure{\(a_i\)}} $b_i\in\Gamma(X',\mathcal O_{X'})$ est la section déduite de $T_i$, alors \inserted{les} $\deleted{\unsure{D}}\Delta_i=V(b_i)$ sont des diviseurs réguliers \unsure{de|dans} $X'$ tels que: $\Delta=\sum\Delta_i$ \unsure{soit|est} un diviseur à croisements normaux (cf. SGA 1 $\overline{\underline{\mathrm{XII}}}$ \unsure{\ldots?}) \inserted{Donc on peut appliquer 4.5. à $X'$ et $U'=X'-\operatorname{Supp}\Delta$.} D'ailleurs, il est immédiat que $\operatorname{Supp}\Delta_i$ n'est autre que l'image inverse de $\operatorname{Supp}D_i$, plus précisément, si $p:X'\unsure{{-}|\longrightarrow}X$ est la projection, on a

\[
p^*(D_i)=m\Delta_i.
\]

On en conclut en particulier

\[
U'=p^{-1}(U).
\]

Considérons l'hom. \deleted{\unsure{can.}} de changement de base

\[
p^*(R^*f_*(\Lambda_U))\longrightarrow R^*f'_*(\Lambda_{U'}),
\tag*{(4.\unsure{6|5}.1)}
\]

où $f':U'-X'$ est l'inclusion. \deleted{\unsure{On}} On constate \inserted{\unsure{alors|donc}} aussitôt sur la définition des isomorphismes (4.5.1) que l'on a commutativité dans le diagramme
\scanpage{115}
\mspage{30}
% archive sheet number: 114
% Upper diagram. The two left vertical signs are unheaded isomorphism signs.
\diagram{
\begin{tikzcd}[ampersand replacement=\&,column sep=large,row sep=large]
\& p^*(R^1f_*(\Lambda_U)) \arrow[r,"(4.6.1)"] \arrow[d,phantom,"{\wr\;\deleted{\unsure{\leftarrow\!\downarrow\!\nearrow|\leftrightarrow\!\downarrow}}}" description] \& R^{\unsure{*|1|\alpha}}f'_*(\Lambda_{U'}) \arrow[dd,"\wr"] \\
(4.6.2) \& p^*\left(\coprod_i\Lambda_{D_i}(-1)\right) \arrow[d,phantom,"\wr" description] \& \\
\& \coprod_i\Lambda_{\Delta_i}(-1) \arrow[r,"m.\mathrm{id}"] \& \coprod_i\Lambda_{\Delta_i}(-1),
\end{tikzcd}
}

où les flèches verticales \deleted{\unsure{de la}} correspondent aux isom. (4.5.1). Lorsque $X$ et $X'$ satisfont au th. de pureté, on en conclut de $\unsure{\overline{m}|\overline{\sim}}$ commutativité dans

% Lower diagram: the long right vertical stroke has no visible arrowhead.
\diagram{
\begin{tikzcd}[ampersand replacement=\&,column sep=large,row sep=large]
(4.6.3) \& p^*(R^pf_*(\Lambda_U)) \arrow[r,"(4.6.1)"] \arrow[d,phantom,"\wr" description] \& R^pf'_*(\Lambda_{U'}) \arrow[dd,dash,"\wr"] \\
\& p^*\bigwedge^p\left(\coprod_i\Lambda_{D_i}(-1)\right) \arrow[d,phantom,"\wr" description] \& \\
\& \bigwedge^p\coprod_i\Lambda_{\Delta_i}(-1) \arrow[r,"m^p.\mathrm{id}"] \& \bigwedge^p\coprod_i\Lambda_{\Delta_i}(-1).
\end{tikzcd}
}

On en conclut en particulier:
